\documentclass{article}%
\usepackage[T1]{fontenc}%
\usepackage[utf8]{inputenc}%
\usepackage{lmodern}%
\usepackage{textcomp}%
\usepackage{lastpage}%
\usepackage{geometry}%
\geometry{margin=2cm}%
\usepackage{expex}%
\usepackage[T1]{fontenc}%
\usepackage[utf8]{inputenc}%
\usepackage[spanish]{babel}%
%
\title{Análisis Lingüístico: Gramatica Normativa Kaqchikel}%
\author{Generado por el TFM de Elías Carrasco}%
%
\begin{document}%
\normalsize%
\maketitle%
\section{Page 89}%
\label{sec:Page89}%
\subsection{Tajb’en techlal amb’il ojetaq moqa otaq / Uso de ojetaq u otaq}%
\label{subsec:Tajbentechlalambilojetaqmoqaotaq/Usodeojetaquotaq}%

%
Se utiliza para indicar que una acción se había terminado antes %
de la ocurrencia de otra acción pero antes del día en que se %
está hablando. %
Ejemplo %
\par\medskip%
\ex
\textit{Ya había venido cuando empezó a llover ayer.} \\
\begingl
\gla Ojetaq chin ule tej ttzaj jb’al. //
\glb  //
\glc  //
\endgl
\xe
%
\medskip%
\par\medskip%
\ex
\textit{Ya habíamos venido cuando veniste.} \\
\begingl
\gla Otaq qo wane tej tula. //
\glb  //
\glc  //
\endgl
\xe
%
\medskip%
\par\medskip%
\ex
\textit{Ya había germinado mi milpa cuando llegué.} \\
\begingl
\gla Ojetaq jawul twi’ nkojine tej npone. //
\glb  //
\glc  //
\endgl
\xe
%
\medskip%
\par\medskip%
\ex
\textit{Cuando viniste ya había leido el libro.} \\
\begingl
\gla Otaq b’aj nsch’ina u’j tej s-ula. //
\glb  //
\glc  //
\endgl
\xe
%
\medskip%
\subsection{Yeky’il tumel / Movimiento y dirección}%
\label{subsec:Yekyiltumel/Movimientoydireccin}%

%
Tienen el papel de indicar la dirección y trayectoria de determinada acción que le describe el hablante al oyente. Todos los direccionales se derivan de los llamados verbos intransitivos de movimiento. Los direccionales pierden parte de sus elementos y son objeto de cambios en su estructura fonológica cuando se combinan con otros direccionales o con verbos transitivos. (Variación dialectal Mam 2000:110) %
K’ulb’il Yol Mam te K’ulb’il Yol Twitz Paxil 89 

%
\end{document}